\section{Discussion and limitations}\label{sec:discussion}

These results are consistent with a simple interpretation. A well-placed early decision changes the state that later rounds inherit. FRO-2 is conservative in the strict sense: it makes fewer first-round flips, far fewer of them wrong, and leaves more true errors to the later rounds. The lighter, cleaner first round is the benign direction; the larger residual is deferred work, not damage. Early error reduction is therefore not a sufficient objective; the sign, timing and size of a displacement jointly govern the outcome.

The specification objection deserves a fair statement. Production BIKE-Flip sustains its margin in every iteration, while our intervention touches a single round; the two claims concern different schedules. The specification calibrates a whole schedule, whereas our destructive case is a mid-flight displacement of a per-round-calibrated one. FRO-2 is not the specification's own first-round protection; it is a one-shot probe on the plain BF decoder.

A randomised control completes the argument by removing position from the schedule. The control runs on BF30, a fixed-threshold BF variant, and adds coin-driven jitter to the threshold in every iteration. Symmetric jitter moves the failure rate from 4.29\% to 19.67\%; a one-sided positive-only recipe lands between 50.5\% and 51.8\%. The same recipes on the plain BF decoder leave no resolved difference (Fisher $p\ge0.54$). Unstructured mixing is not a substitute for a positioned deterministic schedule.

The study has four boundaries. All primary conclusions rest on rates between roughly $10^{-4}$ and $10^{-3}$, so no production-target extrapolation is attempted. Two operating points carry the primary inference and a third is supplemental; other block sizes have behaviour screens only. The offset grid and screened schedules are our construction, and the cross-implementation check is same-project rather than independent. The diagnostic labels use the true error support, which no deployer knows. At production parameters the raised-weight window sits 28\% above the production weight of 134; production-weight evidence is a decision margin and a convergence cost, with a zero-failure bound of $9.4\times10^{-5}$. Capture was not demonstrated; the overlaps are stress constructions. Future work should push measurement toward the production tail and test independent implementations over larger key pools.

\section{Conclusion}\label{sec:conclusion}

The first round of the plain BF decoder is a sharp, asymmetric two-sided window. Conservative offsets help, while the permissive side collapses within two counts. The round and persistence of the offset, not its size alone, govern conditional failure risk. The window is consistent with asymmetric false-flip suppression: the intervention removes wrong decisions rather than accelerating correction. It transfers across points and weights without retuning, stops at the tested capture-prone constructions, and reproduces at the production block size once the error weight is raised; at the production weight only a decision margin and a convergence cost are measurable. A round-one surrogate did not predict another schedule's effect, and randomisation did not reproduce a positioned decision. Practically, pair instances variant by variant, cluster inference by key, and calibrate thresholds where failures are measurable.
